
% =============================================================================
% Project : THE PHILOSOPHY OF VIABILITY
% Book    : Book I — Ontology of Viability
% Book ID : VS.B01
% Part    : Part I — Identity and Manifestation
% Part ID : VS.B01.P01
% Chapter : Chapter 08 — Comparative Philosophical Examination
% Chapter ID : VS.B01.P01.C08
%
% File :
%   2_MAINMATTER/PART_I_IDENTITY_AND_MANIFESTATION/
%   CHAPTER_08_COMPARATIVE_PHILOSOPHICAL_EXAMINATION/
%   sections/section_08_locke.tex
%
% Section :
%   Section 08 — John Locke
%
% Structural Identifier :
%   VS.B01.P01.C08.S08
%
% Purpose :
%   To reconstruct Locke's account of identity across kinds,
%   his distinctions among masses, organisms, human beings,
%   and persons, and his account of personal identity through
%   consciousness; then to examine the explanatory reach and
%   limitations of these distinctions from the ontological
%   commitments and arguments established in Part I.
%
% Scope :
%   The examination addresses identity relative to kinds,
%   the identity of masses and living organisms, the distinction
%   between a human being and a person, consciousness and
%   personal identity, memory, responsibility, and the limits
%   of Locke's account as a general ontology of identity.
%
% Philosophical Restriction :
%   The critical examination is grounded in Part I of Book I.
%   Locke's account of personal identity is treated as a
%   historically specific theory and not as a foundational
%   account of Being or identity in general. No epistemological
%   framework from another Book is imported into the ontology.
%
% Claim Status :
%   The reconstruction of Locke's position is textually
%   supported by his Essay, Book II, Chapter XXVII.
%   The comparative assessment is argued from distinctions
%   established in Part I.
%   The question whether kind-relative criteria can contribute
%   to a general ontological account of identity remains open.
%
% =============================================================================


\section{John Locke}
\label{sec:ch08_s08}

\index[name]{Locke, John}
\index[concept]{Identity}
\index[concept]{Personal Identity}
\index[concept]{Consciousness}
\index[concept]{Memory}
\index[concept]{Man}
\index[concept]{Person}
\index[concept]{Organism}
\index[concept]{Living Organism}
\index[concept]{Mass of Matter}
\index[concept]{Numerical Identity}
\index[concept]{Persistence}
\index[concept]{Continuity}
\index[concept]{Responsibility}
\index[concept]{Kind-Relative Identity Criteria}


% =============================================================================
% SCOPE AND COMPARATIVE ORIENTATION
% =============================================================================

\subsection*{Scope and Comparative Orientation}

% Structural ID: VS.B01.P01.C08.S08.SS01.P001
\phantomsection\label{para:VS.B01.P01.C08.S08.SS01.P001}

John Locke occupies a central position in the history of
philosophical discussions of identity because he explicitly
distinguishes several questions that are often treated as
though they were one. The identity of a mass of matter, the
identity of a living organism, the identity of a human being,
and the identity of a person are not governed, in his account,
by one undifferentiated criterion. His discussion therefore
provides an important comparative test for any ontology that
attempts to clarify what identity means without assuming that
every existent must be individuated or tracked in the same
manner. \cite[II.xxvii]{LockeEssay}

% Structural ID: VS.B01.P01.C08.S08.SS01.P002
\phantomsection\label{para:VS.B01.P01.C08.S08.SS01.P002}

The principal question of this examination is not whether
Locke has discovered a universally applicable test by which
every identity can be recognized. It is whether his
distinction among kinds of things reveals something
ontologically significant about identity, or whether it
primarily establishes different criteria for determining
when an entity of a particular kind is to be counted as the
same entity across a change or an interval of time. These
possibilities must not be conflated. A criterion for judging
sameness may be relevant to identity without constituting
an exhaustive explanation of what identity is.

% Structural ID: VS.B01.P01.C08.S08.SS01.P003
\phantomsection\label{para:VS.B01.P01.C08.S08.SS01.P003}

The examination proceeds from Locke's own distinctions
rather than from a generalized characterization of his
philosophy. It reconstructs his treatment of masses,
organisms, human beings, and persons; considers the role
assigned to consciousness in personal identity; and then
tests the scope of his account against the distinctions
between Being and identity, identity and determination,
identity and individuation, and identity and persistence
developed in Part I. The aim is neither to dismiss
kind-relative criteria nor to treat them as sufficient
without further argument.

% Structural ID: VS.B01.P01.C08.S08.SS01.P004
\phantomsection\label{para:VS.B01.P01.C08.S08.SS01.P004}

A methodological restriction is essential. Locke's account
of personal identity concerns the conditions under which
actions, experiences, and responsibility may be attributed
to a person. That subject is philosophically important, but
it cannot silently replace the broader ontological question
of identity. The present section therefore distinguishes
Locke's theory of personal identity from the general problem
of the identity of beings. The former is examined as one
specific account within the latter problem, not as its
universal foundation.


% =============================================================================
% 1. HISTORICAL AND TEXTUAL LOCATION
% =============================================================================

\subsection{Historical and Textual Location}
\label{subsec:ch08-locke-historical}

% Structural ID: VS.B01.P01.C08.S08.SS02.P001
\phantomsection\label{para:VS.B01.P01.C08.S08.SS02.P001}

John Locke (1632--1704) develops his principal discussion
of identity in Book II, Chapter XXVII of
\textit{An Essay Concerning Human Understanding}. The
chapter is distinctive because it does not begin by
assuming that identity has one uniform criterion applicable
to every subject. Locke instead examines what it means for
different kinds of things to remain the same, and he
distinguishes the identity of a mass of matter from that
of a living organism, a human being, and a person.
\cite[II.xxvii]{LockeEssay}

% Structural ID: VS.B01.P01.C08.S08.SS02.P002
\phantomsection\label{para:VS.B01.P01.C08.S08.SS02.P002}

The chapter's argument is organized around the relation
between identity and the kind of thing whose identity is
under consideration. Locke's examples are not merely
illustrations appended to a prior abstract definition.
They perform a central argumentative function: they show
that the conditions under which something is called the
same depend upon what sort of thing it is taken to be.
An account adequate for a mass of matter need not be
adequate for an organized living body, and neither
criterion can simply be transferred to the identity of
a person.

% Structural ID: VS.B01.P01.C08.S08.SS02.P003
\phantomsection\label{para:VS.B01.P01.C08.S08.SS02.P003}

The distinction is important for the comparative inquiry
because it prevents a premature generalization from one
domain of examples. Locke's argument requires the reader
to ask what kind of entity is being discussed before
applying a criterion of sameness. The ontological force
of this requirement, however, remains to be assessed.
The fact that different kinds require different criteria
does not by itself establish whether those criteria
constitute identity, indicate identity, or specify the
conditions under which identity is attributed.


% =============================================================================
% 2. IDENTITY AND KIND-RELATIVE CRITERIA
% =============================================================================

\subsection{Identity and the Kind of Thing}
\label{subsec:ch08-locke-kinds}

% Structural ID: VS.B01.P01.C08.S08.SS03.P001
\phantomsection\label{para:VS.B01.P01.C08.S08.SS03.P001}

Locke's general approach begins from the proposition
that the identity of a thing must be considered in
relation to the kind under which it is understood.
To ask whether something is the same is incomplete
unless the subject of the question is specified.
The same mass, the same plant, the same animal, the
same human being, and the same person do not necessarily
remain the same according to an identical criterion.
\cite[II.xxvii.1--8]{LockeEssay}

% Structural ID: VS.B01.P01.C08.S08.SS03.P002
\phantomsection\label{para:VS.B01.P01.C08.S08.SS03.P002}

This approach has a significant comparative strength.
It resists the assumption that all entities are
appropriately treated as aggregates of unchanging
material parts. It also resists the assumption that
the persistence of a particular arrangement or
collection of matter must be the criterion of identity
for every kind of existent. The identity of an organism,
for example, cannot be settled merely by listing the
material particles that compose it at one moment and
then comparing that list with the particles present
at another moment.

% Structural ID: VS.B01.P01.C08.S08.SS03.P003
\phantomsection\label{para:VS.B01.P01.C08.S08.SS03.P003}

Nevertheless, the phrase "same kind" requires care.
A kind-relative criterion may specify what counts as
the continuation of an entity under a particular
classification. It does not automatically establish
that the classification itself provides the complete
ontological ground of the entity's identity. Nor does
the dependence of a criterion upon a kind show that
identity is reducible to the definition, description,
or conceptual classification of that kind.

% Structural ID: VS.B01.P01.C08.S08.SS03.P004
\phantomsection\label{para:VS.B01.P01.C08.S08.SS03.P004}

The distinction between a criterion and that which the
criterion is intended to assess is therefore decisive.
A criterion may be necessary for making a determinate
judgment about sameness. Yet its applicability does
not by itself demonstrate that the criterion is
identical with the ontological condition of identity.
The comparative question is whether Locke sufficiently
distinguishes these levels or whether his account
sometimes moves from the conditions of identifying
a thing to the conditions constituting its identity.

% Structural ID: VS.B01.P01.C08.S08.SS03.P005
\phantomsection\label{para:VS.B01.P01.C08.S08.SS03.P005}

From the standpoint of Part I, this question must not
be answered by rejecting kind-relative criteria as
such. The distinctions established in the discussion
of identity and determination, and of identity and
individuation, permit a more exact assessment. A kind
may constrain the relevant determinations and the
appropriate mode of individuation without thereby
becoming identical with identity itself. The criterion
may be domain-sensitive while the ontological question
remains distinct. \SectionRef{sec:identity-and-determination}
\SectionRef{sec:identity-and-individuation}


% =============================================================================
% 3. THE IDENTITY OF MASSES OF MATTER
% =============================================================================

\subsection{Masses of Matter and Material Composition}
\label{subsec:ch08-locke-masses}

% Structural ID: VS.B01.P01.C08.S08.SS04.P001
\phantomsection\label{para:VS.B01.P01.C08.S08.SS04.P001}

Locke distinguishes the identity of a mass of matter
from the identity of a living organism. A mass is
considered in terms of the particles that compose it.
If particles are added to or removed from a particular
mass, the mass is no longer composed of precisely the
same material collection. Locke's discussion therefore
connects the identity of a mass with the continued
presence of the same constituent particles.
\cite[II.xxvii.3--4]{LockeEssay}

% Structural ID: VS.B01.P01.C08.S08.SS04.P002
\phantomsection\label{para:VS.B01.P01.C08.S08.SS04.P002}

This account is intelligible when the subject is
specifically a mass considered as a collection of
material parts. It makes clear why the addition or
subtraction of particles can matter to the sameness
of that collection. Locke thereby identifies a
criterion appropriate to one particular mode of
considering a material entity. His account is not
simply an assertion that every existent must retain
all of its original matter in order to remain itself.

% Structural ID: VS.B01.P01.C08.S08.SS04.P003
\phantomsection\label{para:VS.B01.P01.C08.S08.SS04.P003}

The limitation appears when the criterion is extended
beyond its appropriate subject. If the material
composition of a mass is treated as the universal
criterion of identity, the identity of an organism
becomes difficult to explain. Living organisms undergo
material replacement while continuing to be treated
as the same living beings. The difference between
these cases is not an incidental complication; it
is precisely why Locke distinguishes the identity
of a mass from that of an organism.

% Structural ID: VS.B01.P01.C08.S08.SS04.P004
\phantomsection\label{para:VS.B01.P01.C08.S08.SS04.P004}

The comparative significance is twofold. First,
material constitution cannot be assumed to be the
universal ground of identity. Second, the failure
of material sameness to provide a universal criterion
does not establish that material constitution is
irrelevant in every case. It remains a possible
determination of an entity and may be decisive for
the identity conditions of certain kinds of material
collections. The proper conclusion is therefore
restricted: material composition cannot be generalized
into a sufficient account of identity as such.

% Structural ID: VS.B01.P01.C08.S08.SS04.P005
\phantomsection\label{para:VS.B01.P01.C08.S08.SS04.P005}

This restriction is consistent with the distinction
between identity and composition developed in Part I.
An entity's composition may change, and the significance
of that change may vary according to the kind of entity
under consideration. But composition and identity
must not be equated. Conversely, the fact that
composition is not identical with identity does not
mean that composition can never contribute to the
conditions under which an entity is individuated or
identified. \SectionRef{sec:plurality-and-composition}


% =============================================================================
% 4. THE IDENTITY OF LIVING ORGANISMS
% =============================================================================

\subsection{Living Organisms and the Continuity of Life}
\label{subsec:ch08-locke-organisms}

% Structural ID: VS.B01.P01.C08.S08.SS05.P001
\phantomsection\label{para:VS.B01.P01.C08.S08.SS05.P001}

Locke's account of a living organism differs from
his account of a mass of matter. The identity of a
plant or animal does not consist in the uninterrupted
retention of precisely the same particles. Rather,
the organism is considered through the continuation
of an organized life. The material parts may change
while the living organization continues as the life
of the same organism. \cite[II.xxvii.4--6]{LockeEssay}

% Structural ID: VS.B01.P01.C08.S08.SS05.P002
\phantomsection\label{para:VS.B01.P01.C08.S08.SS05.P002}

The example of the oak is particularly important.
The tree remains the same living organism despite
the replacement of its constituent matter through
growth and other material processes. Its identity
is not therefore reducible to the numerical sameness
of its particles. Locke's account shifts attention
from the continued possession of a fixed collection
of matter to the continuation of an organized life
in a succession of changing material parts.
\cite[II.xxvii.4]{LockeEssay}

% Structural ID: VS.B01.P01.C08.S08.SS05.P003
\phantomsection\label{para:VS.B01.P01.C08.S08.SS05.P003}

This distinction provides an important challenge
to any account that identifies persistence with
the preservation of every constituent. The continued
existence of an organism cannot be explained simply
by asserting that its original material composition
remains unchanged. The organism's material parts
are subject to replacement, while the living being
is nevertheless treated as continuous through that
replacement. Locke's analysis thus makes the relation
between identity and material change a central
ontological problem.

% Structural ID: VS.B01.P01.C08.S08.SS05.P004
\phantomsection\label{para:VS.B01.P01.C08.S08.SS05.P004}

Yet the appeal to the continuation of life must
itself be examined. It identifies a relevant
criterion for living beings, but the expression
"the same life" requires clarification if it is
to provide more than a restatement of the identity
that it is intended to explain. One must specify
what continuity of life entails, which transformations
are compatible with it, and what limits distinguish
the continuation of one organism from the cessation
of one organism and the emergence of another.

% Structural ID: VS.B01.P01.C08.S08.SS05.P005
\phantomsection\label{para:VS.B01.P01.C08.S08.SS05.P005}

This is not a claim that Locke's account is necessarily
circular. It is a demand for an explicit account of
the relation between the criterion and the phenomenon
being explained. If continuity of life is independently
specified, it may provide a substantive criterion.
If it is defined only as whatever makes the organism
the same organism, the explanation risks becoming
circular. The distinction between identity and
persistence developed in Part I makes this problem
visible without requiring the rejection of continuity
as a relevant condition. \SectionRef{sec:identity-and-persistence}

% Structural ID: VS.B01.P01.C08.S08.SS05.P006
\phantomsection\label{para:VS.B01.P01.C08.S08.SS05.P006}

Locke's treatment of organisms therefore has a
genuine explanatory advantage over a purely
compositional account. It recognizes that the identity
conditions of a living being cannot be derived merely
from the persistence of its material components.
Its remaining limitation is that the relation between
organized life, persistence, and identity requires
further specification. The account identifies an
important domain-specific criterion, but it does
not by itself settle the general ontology of identity.


% =============================================================================
% 5. THE IDENTITY OF A HUMAN BEING
% =============================================================================

\subsection{The Identity of a Human Being}
\label{subsec:ch08-locke-human-being}

% Structural ID: VS.B01.P01.C08.S08.SS06.P001
\phantomsection\label{para:VS.B01.P01.C08.S08.SS06.P001}

Locke distinguishes the identity of a human being
from the identity of a person. A human being, in
the relevant discussion, is considered as a living
organism of a particular kind. Its identity is
therefore connected with the continuation of the
same organized life rather than with the uninterrupted
presence of precisely the same material particles
or with the continued existence of an unchanged
substance considered independently of that life.
\cite[II.xxvii.6--8]{LockeEssay}

% Structural ID: VS.B01.P01.C08.S08.SS06.P002
\phantomsection\label{para:VS.B01.P01.C08.S08.SS06.P002}

The distinction prevents the identity of a human
being from being immediately equated with the
identity of a person. A human organism and a person
are not interchangeable terms in Locke's argument.
The former concerns the identity of a living being;
the latter introduces a further account involving
consciousness and the attribution of experiences
and actions. The distinction is central to the
structure of the chapter and must be preserved
when reconstructing Locke's position.

% Structural ID: VS.B01.P01.C08.S08.SS06.P003
\phantomsection\label{para:VS.B01.P01.C08.S08.SS06.P003}

This distinction also exposes a broader conceptual
problem. The identity of an entity under one
description or kind need not settle every question
that can be asked about that entity under another
description. The same living human being may be
considered as an organism and as a person, but
the criteria relevant to these descriptions need
not coincide. Locke's analysis is valuable because
it refuses to treat these questions as automatically
equivalent.

% Structural ID: VS.B01.P01.C08.S08.SS06.P004
\phantomsection\label{para:VS.B01.P01.C08.S08.SS06.P004}

The ontological limitation appears if the distinction
between a human being and a person is taken to imply
that one of these identities is the ultimate or
exclusive identity of the existent. Locke's distinction
establishes that different questions and criteria
are involved. It does not, without additional
argument, establish that the identity of the human
organism is ontologically secondary to personal
identity, or that personal identity exhausts the
identity of the human being.

% Structural ID: VS.B01.P01.C08.S08.SS06.P005
\phantomsection\label{para:VS.B01.P01.C08.S08.SS06.P005}

The distinction between analytical separation and
ontological separation is therefore important.
The human being and the person can be distinguished
for the purposes of Locke's analysis without
immediately establishing that they are two
independent beings or that one is reducible to
the other. An analytical distinction does not
by itself establish an ontological division.
Likewise, the fact that the two descriptions concern
one living human existent does not erase the
conceptual distinction between them. This assessment
preserves the separate roles of Analytical Distinction
and Ontological Inseparability established in Part I.
\SectionRef{sec:being-and-identity}


% =============================================================================
% 6. PERSONAL IDENTITY AND CONSCIOUSNESS
% =============================================================================

\subsection{Personal Identity and Consciousness}
\label{subsec:ch08-locke-person-consciousness}

% Structural ID: VS.B01.P01.C08.S08.SS07.P001
\phantomsection\label{para:VS.B01.P01.C08.S08.SS07.P001}

Locke's account of personal identity introduces a
criterion different from the material composition
of a mass and the continuation of the life of an
organism. He characterizes a person as a thinking,
intelligent being capable of reason and reflection,
and he connects personal identity with the extent
to which consciousness can extend backwards to
past thoughts and actions. The identity of a person
is therefore examined through consciousness rather
than simply through the persistence of a particular
body or the continued existence of a particular
human organism. \cite[II.xxvii.9--11]{LockeEssay}

% Structural ID: VS.B01.P01.C08.S08.SS07.P002
\phantomsection\label{para:VS.B01.P01.C08.S08.SS07.P002}

The role of consciousness in this account is not
merely that of a passive record of events. Locke
connects consciousness with the appropriation of
past experiences and actions as one's own. The
person's relation to a past action is therefore
not explained solely by the fact that the action
occurred in the history of a particular body.
The question is whether the present consciousness
extends to that past action in the relevant sense.
\cite[II.xxvii.9--10]{LockeEssay}

% Structural ID: VS.B01.P01.C08.S08.SS07.P003
\phantomsection\label{para:VS.B01.P01.C08.S08.SS07.P003}

This position gives personal identity a distinctive
practical and forensic significance. Questions
about responsibility, accountability, and the
attribution of actions cannot be treated as though
they were identical with questions about the
persistence of material particles. Locke's account
attempts to explain why a person is appropriately
considered the subject of particular past actions
through the relation of consciousness to those
actions. The forensic dimension is therefore an
integral part of the theory, not an incidental
application added after the account of identity
has been completed.

% Structural ID: VS.B01.P01.C08.S08.SS07.P004
\phantomsection\label{para:VS.B01.P01.C08.S08.SS07.P004}

The strength of Locke's proposal lies in its
recognition that the identity of a person, as he
defines that subject, cannot simply be inferred
from the identity of a mass of matter or from
the uninterrupted identity of a human organism.
It introduces a criterion responsive to the
specific subject under consideration. The account
thereby prevents the question of personal identity
from being collapsed into a purely material
description of the human body.

% Structural ID: VS.B01.P01.C08.S08.SS07.P005
\phantomsection\label{para:VS.B01.P01.C08.S08.SS07.P005}

Nevertheless, the distinction between personal
identity and the identity of the human being does
not establish that consciousness constitutes
identity in the general ontological sense. It
establishes, at most, that consciousness has a
central role in Locke's account of the identity
of persons. The scope of that claim must remain
restricted to the subject Locke is discussing.
A criterion relevant to the identity of persons
cannot be transferred without argument to plants,
animals, material masses, or beings that do not
fall under Locke's definition of a person.

% Structural ID: VS.B01.P01.C08.S08.SS07.P006
\phantomsection\label{para:VS.B01.P01.C08.S08.SS07.P006}

A further distinction is required between consciousness
as a condition of personal identity and consciousness
as a means by which personal identity is recognized
or attributed. These are not necessarily the same
claim. If consciousness is constitutive, Locke must
explain how its relevant extension establishes the
identity of the person. If consciousness functions
as a criterion of attribution, the account must
explain whether it identifies an identity already
obtaining or provides the condition under which
personal identity is assigned. The distinction
between identity and its determination makes this
question unavoidable.

% Structural ID: VS.B01.P01.C08.S08.SS07.P007
\phantomsection\label{para:VS.B01.P01.C08.S08.SS07.P007}

The present comparison does not settle the
interpretation of every aspect of Locke's theory.
It identifies a precise burden of argument: the
relation between consciousness and personal identity
must be specified without assuming that the criterion
and the identity being assessed are identical by
definition. Locke's account is strongest when read
as an attempt to explain the identity conditions
of persons; it becomes more difficult to defend
if it is expanded into a universal ontology of
identity.


% =============================================================================
% 7. MEMORY, CONSCIOUSNESS, AND THE EXTENSION OF IDENTITY
% =============================================================================

\subsection{Memory and the Extension of Consciousness}
\label{subsec:ch08-locke-memory}

% Structural ID: VS.B01.P01.C08.S08.SS08.P001
\phantomsection\label{para:VS.B01.P01.C08.S08.SS08.P001}

Memory is closely connected with Locke's discussion
because it provides an important means through
which consciousness may extend to past experiences.
However, the account should not be simplified into
the claim that personal identity is merely the
possession of a continuous store of memories.
Locke's central formulation concerns consciousness
and its extension, while memory is relevant to
the availability and recovery of past experiences.
The relation between these notions must therefore
be stated carefully. \cite[II.xxvii.9--10, 13--20]{LockeEssay}

% Structural ID: VS.B01.P01.C08.S08.SS08.P002
\phantomsection\label{para:VS.B01.P01.C08.S08.SS08.P002}

The distinction matters because memory can be
incomplete, interrupted, or inaccessible. A person
may fail to remember many events that occurred
during the person's life. If every absent memory
were sufficient to establish the absence of personal
identity, the account would face serious difficulties.
Conversely, the presence of an apparent memory
does not, without further argument, settle every
question about the identity of the subject who
remembers. The relation between memory, consciousness,
and identity cannot be treated as self-evident.

% Structural ID: VS.B01.P01.C08.S08.SS08.P003
\phantomsection\label{para:VS.B01.P01.C08.S08.SS08.P003}

Locke's discussion of forgetting and the limits
of consciousness raises a problem concerning the
scope of the criterion. If personal identity extends
as far as consciousness extends, one must explain
how gaps in consciousness relate to the identity
of a person across time. It is not enough merely
to observe that a person does not presently recall
a particular action. The theory must distinguish
the absence of present recollection from the
absence of the relevant relation of consciousness.

% Structural ID: VS.B01.P01.C08.S08.SS08.P004
\phantomsection\label{para:VS.B01.P01.C08.S08.SS08.P004}

The distinction between memory and identity also
prevents a circular formulation. If a memory is
treated as genuine only when it is already the
memory of the same person who performed the action,
then personal identity is presupposed in the
criterion intended to explain it. If, however,
the memory is treated as sufficient to establish
identity independently, the account must explain
how the relevant relation can be established
without relying upon the identity it is supposed
to determine. This is a substantive difficulty
for the theory, not a reason to dismiss memory
as irrelevant.

% Structural ID: VS.B01.P01.C08.S08.SS08.P005
\phantomsection\label{para:VS.B01.P01.C08.S08.SS08.P005}

The comparative assessment must remain precise.
The difficulty concerns the explanatory relation
between consciousness, memory, and personal identity.
It does not establish that personal identity is
impossible, that consciousness has no relevance
to personal identity, or that all memory-based
criteria are invalid. Nor does it settle the
identity of the human organism. The scope of
the objection is restricted to the sufficiency
of the proposed criterion for explaining the
identity of persons.

% Structural ID: VS.B01.P01.C08.S08.SS08.P006
\phantomsection\label{para:VS.B01.P01.C08.S08.SS08.P006}

From the perspective of Part I, the deeper issue
is whether a relation that connects an existent
with its past manifestations can be treated as
the identity of that existent. Consciousness may
provide a distinctive relation between a person
and past experiences, but a relation to past
manifestations is not automatically identical
with the identity of the subject. The distinction
between identity, continuity, and persistence
must therefore remain in force. \SectionRef{sec:identity-and-persistence}


% =============================================================================
% 8. THE PRINCE AND THE COBBLER
% =============================================================================

\subsection{The Prince and the Cobbler}
\label{subsec:ch08-locke-prince-cobbler}

% Structural ID: VS.B01.P01.C08.S08.SS09.P001
\phantomsection\label{para:VS.B01.P01.C08.S08.SS09.P001}

Locke's example of the prince and the cobbler
is designed to clarify the difference between
the identity of a human being and the identity
of a person. In the imagined case, the consciousness
of the prince is transferred to the body of the
cobbler, while the consciousness associated with
the cobbler's former life is absent. Locke argues
that the person would follow the consciousness
rather than simply the continued identity of the
human body. \cite[II.xxvii.15]{LockeEssay}

% Structural ID: VS.B01.P01.C08.S08.SS09.P002
\phantomsection\label{para:VS.B01.P01.C08.S08.SS09.P002}

The example is not intended merely to describe
a physical transformation. Its argumentative
purpose is to isolate the criterion of personal
identity from the criterion of bodily or organismic
identity. The body associated with one human
being and the consciousness associated with
another are conceptually separated so that the
reader can consider which relation determines
the identity of the person. The example therefore
makes explicit the distinction that structures
Locke's account.

% Structural ID: VS.B01.P01.C08.S08.SS09.P003
\phantomsection\label{para:VS.B01.P01.C08.S08.SS09.P003}

The example has considerable explanatory value
within Locke's theory. It shows that personal
identity, as Locke understands it, is not simply
the identity of a particular human body. It also
demonstrates why the identity of a person cannot
be assumed to coincide with the identity of the
living organism. The case is designed to test
the implications of the consciousness criterion
by separating relations that ordinary experience
usually presents together.

% Structural ID: VS.B01.P01.C08.S08.SS09.P004
\phantomsection\label{para:VS.B01.P01.C08.S08.SS09.P004}

Yet the example does not independently prove
that consciousness constitutes identity. It
demonstrates what follows if Locke's criterion
is accepted and if the imagined separation is
coherent. A thought experiment can clarify the
consequences of a principle without establishing
that the principle is ontologically correct.
The example therefore supports the internal
structure of Locke's argument but cannot, by
itself, settle the general question of identity.

% Structural ID: VS.B01.P01.C08.S08.SS09.P005
\phantomsection\label{para:VS.B01.P01.C08.S08.SS09.P005}

The comparative question is whether the example
shows that one existent can have different
identities under different descriptions, or
whether it distinguishes different subjects
and criteria that should not be collapsed.
The example alone does not settle this broader
question. It establishes the conceptual separation
between human-being identity and personal identity
within Locke's framework, while leaving open
the ontological interpretation of the relation
between them.

% Structural ID: VS.B01.P01.C08.S08.SS09.P006
\phantomsection\label{para:VS.B01.P01.C08.S08.SS09.P006}

The distinction between analytical separation
and ontological priority is particularly important
here. Locke's argument gives personal identity
a criterion distinct from the criterion for a
human organism. It does not follow that personal
identity has ontological priority over organismic
identity, or that the organism is merely an
instrument through which a person manifests.
Such conclusions would require additional premises
that the example itself does not establish.


% =============================================================================
% 9. PERSONAL IDENTITY, RESPONSIBILITY, AND FORENSIC CONSIDERATIONS
% =============================================================================

\subsection{Personal Identity and Responsibility}
\label{subsec:ch08-locke-responsibility}

% Structural ID: VS.B01.P01.C08.S08.SS10.P001
\phantomsection\label{para:VS.B01.P01.C08.S08.SS10.P001}

Locke connects personal identity with the attribution
of actions and responsibility. The person is
considered in relation to the actions and experiences
that can be attributed to that person through
consciousness. This forensic dimension helps
explain why Locke does not simply identify a
person with a body or with an immaterial substance
considered independently of consciousness. The
account is concerned with the subject to whom
actions may properly be attributed. \cite[II.xxvii.16--26]{LockeEssay}

% Structural ID: VS.B01.P01.C08.S08.SS10.P002
\phantomsection\label{para:VS.B01.P01.C08.S08.SS10.P002}

The practical importance of this approach should
not be understated. Questions of responsibility
are not identical with questions about the
persistence of matter. A material body may
continue to exist while the conditions for
attributing a particular action to a person
remain disputed. Locke's account attempts to
make the relation between consciousness and
attribution explicit rather than treating
bodily continuity as sufficient for every
question concerning personal identity.

% Structural ID: VS.B01.P01.C08.S08.SS10.P003
\phantomsection\label{para:VS.B01.P01.C08.S08.SS10.P003}

Nevertheless, the practical function of a
criterion does not by itself establish its
ontological sufficiency. A theory may explain
how responsibility is attributed without
providing a complete account of the identity
of the being to whom responsibility is attributed.
The forensic question and the ontological
question are related, but they are not
interchangeable. The fact that a criterion
serves an important practical purpose does
not prove that it constitutes identity in
every relevant sense.

% Structural ID: VS.B01.P01.C08.S08.SS10.P004
\phantomsection\label{para:VS.B01.P01.C08.S08.SS10.P004}

The distinction also prevents a reverse
reduction. The limitations of a forensic
criterion do not show that responsibility
is philosophically unimportant or that the
identity of persons must be reduced to the
identity of organisms. The appropriate
conclusion is narrower: an account of
responsibility and personal attribution
cannot be treated as a complete substitute
for an ontology of identity unless the
relation between these questions is
independently established.


% =============================================================================
% 10. THE STRONGEST CASE FOR LOCKE
% =============================================================================

\subsection{The Strongest Defence of Locke's Position}
\label{subsec:ch08-locke-defence}

% Structural ID: VS.B01.P01.C08.S08.SS11.P001
\phantomsection\label{para:VS.B01.P01.C08.S08.SS11.P001}

The strongest defence of Locke begins by rejecting
the assumption that a general theory of identity
must provide one identical criterion for every
kind of existent. Such a demand may itself be
misguided. If masses, organisms, and persons
are different kinds of subjects, their identity
conditions may differ without contradiction.
The plurality of criteria may reflect genuine
differences among the entities under consideration
rather than a defect in the theory.

% Structural ID: VS.B01.P01.C08.S08.SS11.P002
\phantomsection\label{para:VS.B01.P01.C08.S08.SS11.P002}

On this defence, Locke's account is not intended
to reduce all identity to consciousness. His
discussion explicitly distinguishes the identity
of masses and organisms from personal identity.
Consciousness is introduced for the specific
case of persons, not as a universal condition
for the identity of everything that exists.
A criticism that attributes to Locke the claim
that plants, animals, and material masses must
all possess consciousness in order to remain
identical would therefore misrepresent his
argument.

% Structural ID: VS.B01.P01.C08.S08.SS11.P003
\phantomsection\label{para:VS.B01.P01.C08.S08.SS11.P003}

The defender may further argue that Locke's
distinction between a human being and a person
is analytically productive. It prevents the
question of personal responsibility from being
settled merely by reference to bodily continuity.
The prince-and-cobbler example is intended to
show that the identity relevant to personal
attribution can diverge from the identity
relevant to the continuation of a human organism.
That divergence is not necessarily a contradiction;
it may disclose different legitimate questions
about one or more subjects.

% Structural ID: VS.B01.P01.C08.S08.SS11.P004
\phantomsection\label{para:VS.B01.P01.C08.S08.SS11.P004}

The defender can also argue that Locke does
not need to provide a complete metaphysical
account of Being in order to make a significant
contribution to the analysis of identity.
A theory may clarify the criteria appropriate
to a particular domain without claiming to
explain every ontological question. If Locke's
purpose is understood as the analysis of
identity across kinds, his account can be
valuable even if it does not establish a
universal ground of identity.

% Structural ID: VS.B01.P01.C08.S08.SS11.P005
\phantomsection\label{para:VS.B01.P01.C08.S08.SS11.P005}

This defence is substantial and must be retained.
The present critique does not refute Locke
merely by observing that his account is not
a complete ontology of identity. A restricted
theory should be evaluated according to the
question it actually addresses. The appropriate
criticism concerns the extent to which the
criteria Locke identifies can be interpreted
as ontological explanations, not whether
those criteria have any explanatory value.


% =============================================================================
% 11. THE STRONGEST OBJECTION FROM PART I
% =============================================================================

\subsection{The Strongest Objection from the Ontology of Viability}
\label{subsec:ch08-locke-objection}

% Structural ID: VS.B01.P01.C08.S08.SS12.P001
\phantomsection\label{para:VS.B01.P01.C08.S08.SS12.P001}

The strongest objection available from Part I is
that Locke's account risks moving between two
different explanatory tasks: specifying the
criterion under which an entity is judged to
remain the same, and explaining the ontological
identity of the entity itself. His distinctions
among masses, organisms, human beings, and
persons establish that different criteria
may be relevant to different subjects. They
do not by themselves establish that these
criteria constitute the identity they are
used to assess.

% Structural ID: VS.B01.P01.C08.S08.SS12.P002
\phantomsection\label{para:VS.B01.P01.C08.S08.SS12.P002}

The objection must not be formulated as the
claim that identity is independent of every
determination or that criteria are irrelevant.
Part I does not require such a conclusion.
Identity and determination are analytically
distinguishable while remaining ontologically
inseparable. Determinations may be relevant
to the way an entity is specified and
distinguished, but identity must not simply
be equated with any one determination or
with a list of determinations. The question
is whether Locke's kind-relative criteria
preserve this distinction.

% Structural ID: VS.B01.P01.C08.S08.SS12.P003
\phantomsection\label{para:VS.B01.P01.C08.S08.SS12.P003}

A related objection concerns individuation.
A criterion that distinguishes one organism
from another, or one person from another,
may explain how entities are counted or
differentiated within a domain. Yet
individuation and identity are not identical.
The fact that a criterion establishes
which entity is being referred to does
not automatically explain what makes that
entity the identity it is. Locke's analysis
is strongest as an account of domain-specific
criteria; its wider ontological significance
requires a further argument. \SectionRef{sec:identity-and-individuation}

% Structural ID: VS.B01.P01.C08.S08.SS12.P004
\phantomsection\label{para:VS.B01.P01.C08.S08.SS12.P004}

A further objection concerns persistence.
The continuation of an organism's life,
the retention of a mass's material parts,
and the extension of consciousness across
past actions are different relations. None
should be treated as automatically identical
with identity itself. The persistence of
a particular feature, material component,
organization, or relation may be relevant
to an identity judgment, but the existence
of such persistence does not alone settle
the ontological question. \SectionRef{sec:identity-and-persistence}

% Structural ID: VS.B01.P01.C08.S08.SS12.P005
\phantomsection\label{para:VS.B01.P01.C08.S08.SS12.P005}

The objection is therefore not that Locke
has selected the wrong single criterion.
His own argument is more differentiated
than such a criticism would suggest.
The problem is whether the theory adequately
explains the relation between identity and
the different criteria it assigns to
different kinds. If those criteria are
treated as sufficient explanations, the
account may reduce identity to material
composition, organized life, or consciousness
in different domains without explaining
the ontological status of identity itself.

% Structural ID: VS.B01.P01.C08.S08.SS12.P006
\phantomsection\label{para:VS.B01.P01.C08.S08.SS12.P006}

The objection also concerns the relation
between plurality and unity. Locke's
distinction among kinds does not require
the conclusion that every kind contains
a separate and unrelated form of identity.
Nor does the unity of an entity require
that its identity be reduced to a single
criterion of composition or continuity.
The possibility of different identity
criteria must be distinguished from the
claim that there are multiple ontologically
unrelated identities. The latter question
remains open and cannot be settled merely
by classifying different kinds of things.
\SectionRef{sec:plurality-of-identities}
\SectionRef{sec:plurality-and-composition}


% =============================================================================
% 12. REPLY, QUALIFICATION, AND LIMITS OF THE CRITIQUE
% =============================================================================

\subsection{Reply and Qualification}
\label{subsec:ch08-locke-reply}

% Structural ID: VS.B01.P01.C08.S08.SS13.P001
\phantomsection\label{para:VS.B01.P01.C08.S08.SS13.P001}

Locke's strongest reply would be that the objection
demands more from his theory than it is intended
to provide. He may maintain that the task is to
identify the conditions under which a thing of
a given kind remains the same, not to produce
an additional entity or property called identity
behind those conditions. If the criteria are
understood as specifying the identity conditions
appropriate to each kind, then the distinction
between identity and its criteria need not imply
that identity is an independently existing
component separate from the entity.

% Structural ID: VS.B01.P01.C08.S08.SS13.P002
\phantomsection\label{para:VS.B01.P01.C08.S08.SS13.P002}

This reply is important because the distinction
between identity and determination must not
be misrepresented. To say that identity is
not reducible to a determination is not to
say that identity exists as an additional
property attached to an otherwise complete
entity. The ontological inseparability of
Being and identity must be preserved alongside
their analytical distinction. A valid critique
of Locke must therefore avoid treating identity
as a separate object that his criteria have
failed to locate.

% Structural ID: VS.B01.P01.C08.S08.SS13.P003
\phantomsection\label{para:VS.B01.P01.C08.S08.SS13.P003}

The reply nevertheless does not remove the
central question. Even if identity conditions
are appropriate to kinds, the account must
clarify what makes a criterion appropriate
to a particular kind and what relation it
bears to the identity of the entity. Otherwise,
the claim that each kind has its own criterion
may describe the distribution of judgments
without explaining the ontological significance
of the distinctions on which those judgments
depend.

% Structural ID: VS.B01.P01.C08.S08.SS13.P004
\phantomsection\label{para:VS.B01.P01.C08.S08.SS13.P004}

The appropriate result is therefore qualified.
Locke succeeds in showing that identity judgments
cannot be responsibly detached from the kinds
of things being considered. He also shows that
material composition, organismic life, and
personal consciousness cannot simply be
substituted for one another. But his account,
as reconstructed here, does not establish
that these criteria provide an exhaustive
general explanation of ontological identity.

% Structural ID: VS.B01.P01.C08.S08.SS13.P005
\phantomsection\label{para:VS.B01.P01.C08.S08.SS13.P005}

This qualification does not establish the
opposite thesis that there must be one
uniform criterion for all identities.
The plurality of criteria may be legitimate,
and the ontological plurality of identities
remains an open problem. The present analysis
therefore rejects both an unwarranted reduction
of identity to one domain-specific criterion
and an unwarranted inference from different
criteria to a plurality of unrelated ontological
principles.


% =============================================================================
% 13. BILATERAL EVALUATION
% =============================================================================

\subsection{Bilateral Evaluation}
\label{subsec:ch08-locke-evaluation}

% Structural ID: VS.B01.P01.C08.S08.SS14.P001
\phantomsection\label{para:VS.B01.P01.C08.S08.SS14.P001}

Locke's principal contribution to the comparative
problem of identity is his insistence that the
identity of a thing must be examined in relation
to the kind of thing under consideration. His
analysis prevents the universalization of the
criterion appropriate to a mass of matter and
distinguishes the continuation of an organism
from the identity of a person. This is a
substantial corrective to accounts that assume
that one form of material or temporal sameness
can explain every case.

% Structural ID: VS.B01.P01.C08.S08.SS14.P002
\phantomsection\label{para:VS.B01.P01.C08.S08.SS14.P002}

His account of personal identity also brings
into focus the relation between consciousness,
past experience, and the attribution of actions.
It demonstrates that questions about persons
cannot simply be collapsed into questions
about the material composition of a body.
The forensic dimension of the theory explains
why the identity of a person is not treated
as interchangeable with the identity of a
human organism.

% Structural ID: VS.B01.P01.C08.S08.SS14.P003
\phantomsection\label{para:VS.B01.P01.C08.S08.SS14.P003}

The principal limitation is the unresolved
relation between a criterion of identity and
the ontological identity that the criterion
is meant to clarify. Locke's account offers
different criteria for different kinds, but
the transition from kind-relative criteria
to a general explanation of identity requires
further argument. His account of consciousness
is especially restricted to personal identity
and cannot be generalized into a universal
principle of Being or identity.

% Structural ID: VS.B01.P01.C08.S08.SS14.P004
\phantomsection\label{para:VS.B01.P01.C08.S08.SS14.P004}

The evaluation must therefore remain bilateral.
Locke should not be criticized for failing
to provide a theory he did not set out to
provide. His differentiated account is
valuable within its intended domain. At the
same time, its explanatory success in
distinguishing identity criteria does not
settle the wider ontological questions
concerning Being, identity, determination,
individuation, composition, and persistence.

% Structural ID: VS.B01.P01.C08.S08.SS14.P005
\phantomsection\label{para:VS.B01.P01.C08.S08.SS14.P005}

The comparative result can be stated as follows:
Locke establishes the importance of kind-relative
identity conditions, but the ontological status
of those conditions remains open. His analysis
is strongest as a differentiated account of
identity criteria and weaker as a complete
general ontology of identity. This conclusion
preserves the genuine insight of his position
without extending it beyond the arguments
that support it.


% =============================================================================
% 14. OPEN QUESTIONS AND TRANSITION
% =============================================================================

\subsection{Open Ontological Questions}
\label{subsec:ch08-locke-open-questions}

% Structural ID: VS.B01.P01.C08.S08.SS15.P001
\phantomsection\label{para:VS.B01.P01.C08.S08.SS15.P001}

Several questions remain unresolved. First, what
is the precise relation between a kind and the
identity conditions appropriate to it? Does
the kind determine which features are relevant
to identity judgments, or does it also constrain
the ontological constitution of the entity
whose identity is at issue? The distinction
between these possibilities requires further
analysis.

% Structural ID: VS.B01.P01.C08.S08.SS15.P002
\phantomsection\label{para:VS.B01.P01.C08.S08.SS15.P002}

Second, can the continuity of an organism's life
be specified independently of the identity
it is intended to explain? If so, which forms
of transformation are compatible with that
continuity, and what limits distinguish the
continuation of one organism from the emergence
of another? The question cannot be answered
merely by appealing to material sameness, since
Locke's account explicitly distinguishes the
organism from a mass of matter.

% Structural ID: VS.B01.P01.C08.S08.SS15.P003
\phantomsection\label{para:VS.B01.P01.C08.S08.SS15.P003}

Third, what is the precise relation between
consciousness and personal identity? Is
consciousness constitutive of personal identity,
a criterion for its attribution, or a relation
through which past actions become attributable
to a person? The different possibilities
have distinct implications and should not
be treated as interchangeable.

% Structural ID: VS.B01.P01.C08.S08.SS15.P004
\phantomsection\label{para:VS.B01.P01.C08.S08.SS15.P004}

Fourth, does the plurality of identity criteria
indicate a plurality of identities, a plurality
of modes of determination, or different
analytical perspectives upon an entity? The
present examination does not decide among
these alternatives. It establishes only that
the distinction among criteria cannot, without
additional premises, settle the ontological
question of the plurality of identities.

% Structural ID: VS.B01.P01.C08.S08.SS15.P005
\phantomsection\label{para:VS.B01.P01.C08.S08.SS15.P005}

Fifth, how should the relation between identity
and persistence be specified when the relevant
continuity differs across masses, organisms,
and persons? The answer must avoid reducing
identity to persistence while also avoiding
the opposite assumption that persistence is
irrelevant to every identity question. A
more complete account must state the relevant
transformations, domains, and limits without
equating continuity with identity.

% Structural ID: VS.B01.P01.C08.S08.SS15.P006
\phantomsection\label{para:VS.B01.P01.C08.S08.SS15.P006}

These questions remain open within the present
comparative examination. They should not be
resolved by introducing an unestablished
principle or by transferring conclusions
from the theory of personal identity to
the ontology of all beings. Their function
is to identify the precise points at which
Locke's analysis contributes to the problem
of identity and where further ontological
argument is still required.


\subsection{Transition to the Next Comparative Examination}
\label{subsec:ch08-locke-transition}

% Structural ID: VS.B01.P01.C08.S08.SS16.P001
\phantomsection\label{para:VS.B01.P01.C08.S08.SS16.P001}

Locke's examination has shown that identity
criteria may differ according to the kind
of entity under consideration. His distinctions
among masses, organisms, human beings, and
persons expose the inadequacy of treating
material composition, organized life, and
consciousness as interchangeable criteria.
At the same time, the analysis has left open
the ontological relation between these
criteria and the identity of the entities
to which they apply.

% Structural ID: VS.B01.P01.C08.S08.SS16.P002
\phantomsection\label{para:VS.B01.P01.C08.S08.SS16.P002}

The next comparative examination must therefore
consider whether a different account of identity
can clarify the relation among continuity,
change, and the identity of an existent without
reducing identity to a single criterion.
The transition does not presuppose that the
questions raised by Locke have been resolved.
It carries forward the distinction between
identity and the criteria through which identity
is assessed, while leaving the plurality of
identities and the ontological status of
persistence open for further examination.


% =============================================================================
% CLAIM-STATUS AND ARGUMENTUAL SUMMARY
% =============================================================================

\subsection*{Claim-Status Summary}

The claims advanced in this section have different
logical forms and different justificatory statuses.
These two dimensions must not be conflated.

\paragraph{Functional or logical forms.}

\begin{itemize}
    \item \textbf{Historical reconstruction:}
    Locke distinguishes the identity of masses,
    organisms, human beings, and persons, and
    assigns consciousness a central role in
    his account of personal identity.

    \item \textbf{Comparative propositions:}
    Material composition, organismic life, and
    consciousness are not interchangeable
    criteria. Their relevance depends upon
    the kind of subject under consideration.

    \item \textbf{Critical arguments:}
    A criterion of identity must not be
    automatically equated with the ontological
    identity it is intended to assess.

    \item \textbf{Open problems:}
    The relation between kind-relative criteria,
    ontological identity, persistence, and the
    possible plurality of identities remains
    unresolved.
\end{itemize}

\paragraph{Epistemic and justificatory statuses.}

\begin{itemize}
    \item \textbf{Textually supported reconstruction:}
    The account of Locke's distinctions is
    grounded in Book II, Chapter XXVII of
    \textit{An Essay Concerning Human Understanding}.

    \item \textbf{Argued or derived claims:}
    The distinction between identity and
    determination, individuation, composition,
    and persistence is applied to assess the
    explanatory scope of Locke's position.

    \item \textbf{Conditionally supported assessment:}
    Locke's kind-relative criteria are
    explanatory within their respective
    domains, provided that the criteria are
    not treated as exhaustive explanations
    of ontological identity.

    \item \textbf{Unresolved questions:}
    Whether kind-relative criteria correspond
    to a plurality of identities, a plurality
    of determinations, or different analytical
    perspectives remains open.
\end{itemize}


% =============================================================================
% END OF SECTION
% =============================================================================